\chapter{Fenwick Trees and Segment Trees}\label{ch:segtree}

When the array changes between queries, prefix sums break: a single update would invalidate $\Oh{n}$ prefix values. Fenwick trees
and segment trees store \emph{partial} aggregates, arranged so that both an update and a query touch only $\Oh{\log n}$ of them.
The Fenwick tree is the smallest and fastest structure for sums; the segment tree is the general one, which adapts to almost any
associative summary, to range updates through lazy tags, to descents that search instead of summing, and to persistence. This chapter
builds both in detail and then turns to the offline sweeps that make many two-dimensional questions one-dimensional.

\begin{chapterintro}{\gold\ \plat\ \ioi}{Chapter~\ref{ch:static} (static range queries) and Chapter~\ref{ch:bits}
(lowest set bits).}
  \item use Fenwick trees for point and range updates, order statistics by descent, and two dimensions;
  \item build segment trees bottom-up, prove the $\Oh{\log n}$ bounds, and design nodes for non-trivial merges;
  \item search a segment tree by descending it, and update ranges with lazy propagation;
  \item compute the area of a union of rectangles, keep every version with persistence, and know when segment tree beats apply;
  \item turn offline two-dimensional conditions into a sweep plus a Fenwick tree.
\end{chapterintro}

% ======================================================================================================================
\section[Fenwick trees]{Fenwick trees \hfill \gold}\label{sec:fenwick}

A \defn{Fenwick tree} (binary indexed tree) over $a_1, \dots, a_n$ stores in cell $i$ the sum of the $\mathit{lowbit}(i)$ elements
that end at position $i$, where $\mathit{lowbit}(i) = i \mathbin{\&} (-i)$ is the value of the lowest set bit of $i$
(Figure~\ref{fig:fenwick}). Cell 12 $= 1100_2$ covers 4 elements, positions 9 to 12; cell 13 covers only itself.

\begin{itemize}
  \item \textbf{Prefix sum up to $p$}: add cell $p$, remove the lowest set bit of $p$, and repeat until $p = 0$. The blocks visited
        are consecutive and tile $[1, p]$ exactly: $13 = 1101_2 \to 1100_2 \to 1000_2$ visits the blocks $[13, 13]$, $[9, 12]$ and
        $[1, 8]$. One step per set bit: at most $\log_2 n$.
  \item \textbf{Adding to position $p$}: every cell whose block contains $p$ must change. Those are $p$, then $p +
        \mathit{lowbit}(p)$, and so on: adding the lowest set bit carries into the next block that covers $p$. Again at most
        $\log_2 n$ steps.
\end{itemize}

\begin{figure}[htbp]
\centering
\begin{tikzpicture}[x=0.86cm]
  \foreach \ii in {1, ..., 16} \node[font=\ttfamily\scriptsize, text=fgblue] at (\ii, 0) {\ii};
  \foreach \ii/\bl/\lv in {1/1/0, 2/2/1, 3/1/0, 4/4/2, 5/1/0, 6/2/1, 7/1/0, 8/8/3, 9/1/0, 10/2/1, 11/1/0, 12/4/2, 13/1/0,
      14/2/1, 15/1/0, 16/16/4}
    \draw[line width=0.7pt, draw=fgline, fill=fgfill, rounded corners=1.5pt] (\ii - \bl + 0.58, 0.4 + 0.42*\lv)
      rectangle (\ii + 0.42, 0.68 + 0.42*\lv);
  \foreach \ii/\bl/\lv in {13/1/0, 12/4/2, 8/8/3}
    \draw[line width=1.2pt, draw=fgorange!90!black, fill=fgorange!25, rounded corners=1.5pt] (\ii - \bl + 0.58, 0.4 + 0.42*\lv)
      rectangle (\ii + 0.42, 0.68 + 0.42*\lv);
  \foreach \ii/\bl/\lv in {5/1/0, 6/2/1, 8/8/3, 16/16/4}
    \draw[line width=1.6pt, draw=fgblue, rounded corners=1.5pt] (\ii - \bl + 0.58, 0.4 + 0.42*\lv) rectangle
      (\ii + 0.42, 0.68 + 0.42*\lv);
  \node[notel, anchor=north, align=center, text width=14cm] at (8.5, -0.3) {Orange: $\mathit{prefix}(13)$ adds the cells 13, 12 and 8
    ($1101_2 \to 1100_2 \to 1000_2$).\\Blue frames: $\mathit{add}(5)$ changes the cells 5, 6, 8 and 16 ($101_2 \to 110_2 \to
    1000_2 \to 10000_2$), every block that contains position 5.};
\end{tikzpicture}
\caption{The blocks of a Fenwick tree with 16 cells: cell $i$ covers the $\mathit{lowbit}(i)$ positions ending at $i$.}
\label{fig:fenwick}
\end{figure}

The template stores the cells in an array of $n + 1$ numbers and uses 0-based positions in its interface; the worked problem Dynamic
Range Sum Queries (Section~\ref{wp:bit}) is the basic use. It needs half the memory of a bottom-up segment tree and its loops are
shorter, so it is the default for sums and counts.

\usetemplate{Fenwick tree}{Ctrl+Alt+T, I}{cpbit}

\paragraph{Range updates, point queries.} To add $u$ to every element of $[a, b]$ and read single elements, keep a Fenwick tree over
the \emph{differences}: add $u$ at position $a$ and $-u$ at $b + 1$. The value of element $k$ is its original value plus the prefix sum of
the differences up to $k$.

\begin{problem}{Range Update Queries (CSES)}
An array receives $q$ operations: add $u$ to every value in $[a, b]$, or report the value at position $k$ ($n, q \le 2 \cdot 10^5$).
\end{problem}

\begin{code}
void solve() {  // BIT (cpbit) is pasted above solve()
    int n, q;
    cin >> n >> q;
    vi x(n);
    for (ll& v : x) cin >> v;
    BIT d(n + 1);  // the differences added so far: x[k] grew by d[0..k]
    while (q--) {
        int type;
        cin >> type;
        if (type == 1) {
            int a, b;
            ll u;
            cin >> a >> b >> u;
            d.add(a - 1, u), d.add(b, -u);  // +u from a on, -u after b
        } else {
            int k;
            cin >> k;
            print(x[k - 1] + d.prefix(k));  // the differences at 0..k-1
        }
    }
}
\end{code}

\paragraph{Order statistics by descent.} If the Fenwick tree holds counts (1 for every element present), ``the $k$-th element'' is
the first position whose prefix sum reaches $k$. Binary search over prefix sums costs $\Oh{\log^2 n}$; descending through the powers of
two costs $\Oh{\log n}$. Start at position 0 and, for $2^j = $ the highest power of two down to 1, look at the cell
$\mathit{pos} + 2^j$: it holds exactly the count of the block just after $\mathit{pos}$, so if that count is below $k$, jump over the
block and subtract it (Figure~\ref{fig:fendescent}).

\begin{figure}[htbp]
\centering
\begin{tikzpicture}
  \node[blbl, anchor=east] at (-0.2, 0) {present};
  \arrs{P}{0}{0}{1/n, 1/n, 0/dim, 1/n, 1/n, 0/dim, 1/hot, 1/n}
  \foreach \ii in {1, ..., 8} \node[lbl] at ($(P-\ii.south)+(0, -0.22)$) {\ii};
  \node[blbl, anchor=east] at (-0.2, -1.1) {cell $t[i]$};
  \arrs{T}{0}{-1.1}{1/n, 2/n, 0/n, 3/good, 1/n, 1/good, 1/cool, 6/cool}
  \node[notel, anchor=west, text width=7.2cm] at (6.2, -0.55) {Find the 5th present element. Cell 8 holds 6, not below 5:
    stay. Cell 4 holds 3 $< 5$: jump to position 4, $k = 2$. Cell $4 + 2 = 6$ holds 1 $< 2$: jump to 6, $k = 1$. Cell
    $6 + 1 = 7$ holds 1, not below 1: stay. The answer is the position after 6, that is 7 (red). Blue: looked at and not
    taken; green: jumped over.};
\end{tikzpicture}
\caption{Fenwick descent: one cell per power of two, $\Oh{\log n}$ for the $k$-th element.}
\label{fig:fendescent}
\end{figure}

\begin{problem}{List Removals (CSES)}
A list of $n \le 2 \cdot 10^5$ numbers loses one element at a time, given by its current position; print the removed elements.
\solution A Fenwick tree with a 1 for every element still in the list. The element at current position $p$ is the $p$-th 1:
\cc{lower_bound(p)} finds it by descent, and removing it subtracts 1.
\end{problem}

\begin{code}
void solve() {  // BIT (cpbit) is pasted above solve()
    int n;
    cin >> n;
    vi x(n);
    for (ll& v : x) cin >> v;
    BIT alive(n);  // a 1 at every position still in the list
    rep(i, n) alive.add(i, 1);
    vi out;
    rep(j, n) {
        int p;
        cin >> p;
        int i = alive.lower_bound(p);  // the p-th remaining position
        out.push_back(x[i]);
        alive.add(i, -1);
    }
    print(out);
}
\end{code}

\begin{trick}[More from one Fenwick tree]
\begin{itemize}
  \item \textbf{Inversions}: scan from left to right and, for each $a_i$, count the earlier elements larger than it with a Fenwick
        tree over the (compressed) values; $\Oh{n \log n}$.
  \item \textbf{An ordered multiset of small values}: counts per value, insertion and deletion are $\pm 1$, ``how many below
        $x$'' is a prefix sum and ``the $k$-th smallest'' a descent. Faster than a balanced tree when the values are at most
        about $10^6$ (Salary Queries, offline, after compressing).
  \item \textbf{Prefix maxima}: when values only increase, a Fenwick tree with $\max$ instead of $+$ answers prefix maxima; this
        is the $\Oh{n \log n}$ LIS with weights.
\end{itemize}
\end{trick}

\subsection{Range updates with range sums}

To add $v$ to a range \emph{and} ask for range sums, keep two Fenwick trees over differences. If $D$ is the difference array,
\[ \sum_{i < p} a_i = \sum_{i < p} \sum_{j \le i} D_j = \sum_{j < p} D_j\,(p - j) = p \sum_{j < p} D_j - \sum_{j < p} j\,D_j, \]
so one tree keeps $D_j$ and the other $j \cdot D_j$, and a prefix sum is $p \cdot \mathit{first}(p) - \mathit{second}(p)$. A range
update changes two entries of each tree.

\usetemplate{Fenwick tree with range updates}{Ctrl+Alt+Y, I}{cpbitrange}

\subsection{Two dimensions}

Nesting the two loops gives a Fenwick tree over a grid: cell $(i, j)$ covers the rows $(i - \mathit{lowbit}(i), i]$ and the columns
$(j - \mathit{lowbit}(j), j]$, and both operations cost $\Oh{\log n \log m}$. The memory is $\Oh{nm}$, so it suits grids up to a few
thousand on a side; for sparse points, use the offline sweep of Section~\ref{sec:sweepqueries}.

\begin{problem}{Forest Queries II (CSES)}
In an $n \times n$ grid ($n \le 1000$) trees appear and disappear one cell at a time; report the number of trees in rectangles
($q \le 2 \cdot 10^5$ operations).
\end{problem}

\begin{code}
void solve() {  // BIT2D (cpbit2d) is pasted above solve()
    int n, q;
    cin >> n >> q;
    vector<string> g(n);
    for (auto& row : g) cin >> row;
    BIT2D b(n, n);
    rep(r, n) rep(c, n) if (g[r][c] == '*') b.add(r, c, 1);
    while (q--) {
        int type;
        cin >> type;
        if (type == 1) {  // a tree appears or disappears
            int y, x;
            cin >> y >> x;
            y--, x--;
            b.add(y, x, g[y][x] == '*' ? -1 : 1);
            g[y][x] = g[y][x] == '*' ? '.' : '*';
        } else {
            int y1, x1, y2, x2;
            cin >> y1 >> x1 >> y2 >> x2;
            print(b.sum(y1 - 1, x1 - 1, y2, x2));
        }
    }
}
\end{code}

\usetemplate{2D Fenwick tree}{Ctrl+Alt+Y, Shift+I}{cpbit2d}

Building the tree with $n^2$ single additions costs $\Oh{n^2 \log^2 n}$, here $10^8$ cheap steps; an $\Oh{n^2}$ build exists (each cell
adds itself to its parent cell, row by row and then column by column), but the whole program already runs in 0.21\,s.

% ======================================================================================================================
\section[Segment trees]{Segment trees \hfill \gold}\label{sec:segtree}

A \defn{segment tree} is a complete binary tree over the array: every node stores the aggregate of a range, the root the whole array,
the leaves single elements, and the children of a node split its range in half (Figure~\ref{fig:segtree}). Any associative operation
works: sums, minima, gcds, matrix products, or custom summaries.

\begin{figure}[htbp]
\centering
\begin{tikzpicture}[s/.style={seg, minimum height=6mm}]
  \node[s, minimum width=11.8cm] at (0, 0) {$[0, 8)$\quad 27};
  \node[s, minimum width=5.8cm] at (-3, -0.85) {$[0, 4)$\quad 13};
  \node[s, minimum width=5.8cm, cool] at (3, -0.85) {$[4, 8)$\quad 14};
  \foreach \px/\tx in {-4.5/{$[0, 2)$\quad 4}, 1.5/{$[4, 6)$\quad 8}, 4.5/{$[6, 8)$\quad 6}}
    \node[s, minimum width=2.8cm] at (\px, -1.7) {\tx};
  \node[s, minimum width=2.8cm, cool] at (-1.5, -1.7) {$[2, 4)$\quad 9};
  \foreach \val [count=\ii from 0] in {3, 1, 4, 5, 2, 6, 5, 1} \node[s, minimum width=1.3cm] (L\ii) at (-5.25 + \ii*1.5, -2.55) {\val};
  \node[s, minimum width=1.3cm, cool] at (L1) {1};
  \node[notel, align=center, text width=12cm] at (0, -3.35) {The query $[1, 8)$ uses the nodes $[1, 2)$, $[2, 4)$ and $[4, 8)$
    (blue): $1 + 9 + 14 = 24$. A node is used when its range lies inside the query and its parent's does not.};
\end{tikzpicture}
\caption{A segment tree of sums over $3, 1, 4, 5, 2, 6, 5, 1$.}
\label{fig:segtree}
\end{figure}

\paragraph{Why $\Oh{\log n}$.} An update changes the leaf and its ancestors, one node per level. For a query, call a node \emph{used}
if its range lies inside $[l, r)$ but its parent's does not. On each level, at most two nodes are used: if three were, the middle one's
sibling would also lie inside the query, so their common parent would lie inside too, and the middle node would not be used. So a
query combines at most $2 \log_2 n$ nodes.

\paragraph{Bottom-up implementation.} The template stores the tree in an array of $2n$ cells: the leaves in cells $n, \dots, 2n - 1$
and the parent of cell $i$ in cell $i / 2$. An update sets a leaf and walks up with $i \mathrel{/}= 2$; a query moves two pointers $l$
and $r$ up from the leaves, taking the cell at $l$ whenever $l$ is a right child and the cell left of $r$ whenever $r$ is a right
child, which are exactly the used nodes of the argument above. No recursion and no power-of-two padding; it is typically about
twice as fast as a recursive tree. For operations that are not commutative,
the template keeps the left and right parts separately so that the order is preserved.

\usetemplate{Segment tree}{Ctrl+Alt+T, S}{cpseg}

% ======================================================================================================================
\section[Descending a segment tree]{Descending a segment tree \hfill \plat}\label{sec:segdescent}

A query asks ``what is the aggregate of this range?''. Many problems ask the inverse: ``where is the first position at which the
aggregate reaches $x$?''. Binary search over queries costs $\Oh{\log^2 n}$; walking down the tree costs $\Oh{\log n}$: at each node,
look at the left child, go left if the answer is there, otherwise go right (and, for prefix aggregates, remember what the left child
contributed).

\begin{problem}{Hotel Queries (CSES)}
Hotels in a row have $h_i$ free rooms. Each of $m$ groups, in order, takes the first hotel with at least $r_j$ free rooms; report the
hotel or 0 ($n, m \le 2 \cdot 10^5$).
\solution A maximum segment tree over the free rooms. If the root's maximum is below $r_j$, no hotel fits; otherwise walk down,
always into the left child when it has enough rooms (Figure~\ref{fig:hotel}), then subtract $r_j$ at the leaf and repair the
path. $\Oh{\log n}$ per group, 0.15\,s.
\end{problem}

\begin{figure}[htbp]
\centering
\begin{tikzpicture}[s/.style={seg, minimum height=6mm}]
  \node[s, minimum width=11.8cm, cool] at (0, 0) {max 6};
  \node[s, minimum width=5.8cm, cool] at (-3, -0.85) {max 4};
  \node[s, minimum width=5.8cm] at (3, -0.85) {max 6};
  \node[s, minimum width=2.8cm, dim] at (-4.5, -1.7) {max 3};
  \node[s, minimum width=2.8cm, cool] at (-1.5, -1.7) {max 4};
  \node[s, minimum width=2.8cm] at (1.5, -1.7) {max 5};
  \node[s, minimum width=2.8cm] at (4.5, -1.7) {max 6};
  \foreach \val/\sy [count=\ii from 0] in {3/n, 2/n, 4/hot, 1/n, 5/n, 5/n, 2/n, 6/n}
    \node[s, minimum width=1.3cm, \sy] at (-5.25 + \ii*1.5, -2.55) {\val};
  \node[notel, align=center, text width=12cm] at (0, -3.35) {A group needs 4 rooms. The root's maximum 6 says that some hotel
    fits. The left child (max 4) can take it: go left. Its left child (max 3, grey) cannot: go right. There the left leaf has 4:
    hotel 3 (red), the leftmost that fits.};
\end{tikzpicture}
\caption{Descending a maximum segment tree to the leftmost leaf with value at least 4.}
\label{fig:hotel}
\end{figure}

\begin{code}
void solve() {
    int n, m;
    cin >> n >> m;
    int size = 1;
    while (size < n) size *= 2;
    vi t(2 * size, 0);  // a max segment tree over the free rooms
    rep(i, n) cin >> t[size + i];
    per(i, 1, size) t[i] = max(t[2 * i], t[2 * i + 1]);
    vi ans;
    rep(j, m) {
        ll need;
        cin >> need;
        if (t[1] < need) {  // no hotel has enough free rooms
            ans.push_back(0);
            continue;
        }
        int v = 1;  // walk down to the leftmost leaf with enough rooms
        while (v < size) v = t[2 * v] >= need ? 2 * v : 2 * v + 1;
        ans.push_back(v - size + 1);
        for (t[v] -= need; v > 1; v /= 2) t[v / 2] = max(t[v], t[v ^ 1]);
    }
    print(ans);
}
\end{code}

The tree here is padded to a power of two, because the descent needs every node's children to be its two halves; the bottom-up
template without padding does not have that shape. The custom-node template offers the general version, \cc{max_right(l, pred)}:
the largest $r$ such that the aggregate of $[l, r)$ still satisfies a monotone predicate.

% ======================================================================================================================
\section[Designing a node]{Designing a node \hfill \plat}\label{sec:segnode}

A segment tree works for any summary that two neighbouring ranges can merge. The question to ask is: \emph{what must the merge of two
halves know about each half?} For the maximum subarray sum, the best subarray of a range lies in the left half, in the right half, or
crosses the middle, and a crossing one is a best suffix of the left half plus a best prefix of the right half. So a node holds four
numbers: the total, the best prefix, the best suffix and the best subarray (Figure~\ref{fig:segnode}), and
\[
  \mathit{best} = \max(\mathit{best}_L,\ \mathit{best}_R,\ \mathit{suf}_L + \mathit{pref}_R), \quad
  \mathit{pref} = \max(\mathit{pref}_L,\ \mathit{sum}_L + \mathit{pref}_R),
\]
and symmetrically for the suffix.

\begin{figure}[htbp]
\centering
\begin{tikzpicture}
  \node[blbl, anchor=east] at (-0.2, 0) {left half};
  \arrs{L}{0}{0}{2/n, -3/n, 4/cool, -1/cool}
  \node[blbl, anchor=east] at (3.9, 0) {};
  \arrs{R}{3.0}{0}{3/cool, 2/cool, -6/n, 1/n}
  \node[blbl, anchor=west] at (5.95, 0) {right half};
  \node[notel, anchor=north west, text width=13cm] at (-1.6, -0.6) {Left: sum 2, best prefix $2 - 3 + 4 = 3$, best suffix $4 - 1 =
    3$, best 4. Right: sum 0, best prefix $3 + 2 = 5$, best suffix 1, best 5. Merged: the best crossing subarray is the left
    suffix plus the right prefix, $3 + 5 = 8$ (blue), larger than 4 and 5; the merged best prefix is $\max(3, 2 + 5) = 7$.};
\end{tikzpicture}
\caption{The four numbers of a maximum-subarray node, and how two nodes merge.}
\label{fig:segnode}
\end{figure}

\begin{problem}{Subarray Sum Queries (CSES)}
An array of $n \le 2 \cdot 10^5$ values receives $m$ point assignments; after each, report the maximum sum of a subarray (the empty
subarray, with sum 0, is allowed).
\end{problem}

\begin{code}
void solve() {  // SegNode (cpsegnode) is pasted above solve()
    int n, m;
    cin >> n >> m;
    vector<Node> leaves(n);
    for (auto& nd : leaves) {
        ll x;
        cin >> x;
        nd = Node::leaf(x);
    }
    SegNode st(leaves);
    while (m--) {
        int k;
        ll x;
        cin >> k >> x;
        st.set(k - 1, Node::leaf(x));
        print(max(0LL, st.query(0, n).best));  // the empty subarray counts
    }
}
\end{code}

\usetemplate{Custom-node segment tree}{Ctrl+Alt+Y, N}{cpsegnode}

\begin{trick}[Nodes for other questions]
\begin{itemize}
  \item \textbf{Maximum prefix sum of a range} (Prefix Sum Queries): sum and best prefix, the first two fields above.
  \item \textbf{Number of distinct adjacent pairs, or of runs}: the count plus the first and last element of the range, so the merge
        can tell whether the boundary creates a new pair.
  \item \textbf{$\min_k (p_k + |k - i|)$} (Pizzeria Queries): split the absolute value at $i$; keep one tree of $p_k - k$ for $k \le i$
        and one of $p_k + k$ for $k \ge i$, and add back $+i$ or $-i$.
  \item \textbf{Is the range a permutation, a palindrome, sorted?}: hashes, minima, or ``sorted and the boundary is in order''.
\end{itemize}
\end{trick}

% ======================================================================================================================
\section[Lazy propagation]{Lazy propagation \hfill \plat}\label{sec:lazy}

A range \emph{update} (add $v$ to every element of $[l, r)$) would touch every leaf below the range. Instead, apply it to the
$\Oh{\log n}$ nodes that cover the range exactly and leave a \emph{tag} on them: ``my children still owe this update''
(Figure~\ref{fig:lazy}). When a later operation needs to go below a tagged node, it first \emph{pushes} the tag to the two children,
updating their aggregates and composing it with their own tags. Every operation still visits $\Oh{\log n}$ nodes.

\begin{figure}[htbp]
\centering
\begin{tikzpicture}[s/.style={seg, minimum height=6mm}]
  \node[s, minimum width=11.8cm] at (0, 0) {$[0, 8)$\quad sum $27 + 30$};
  \node[s, minimum width=5.8cm] at (-3, -0.85) {$[0, 4)$\quad sum $13 + 10$};
  \node[s, minimum width=5.8cm, hot] at (3, -0.85) {$[4, 8)$\quad sum $14 + 20$, tag $+5$};
  \node[s, minimum width=2.8cm] at (-4.5, -1.7) {$[0, 2)$\quad 4};
  \node[s, minimum width=2.8cm, hot] at (-1.5, -1.7) {$[2, 4)$\ \ $9 + 10$, tag $+5$};
  \node[s, minimum width=2.8cm, dim] at (1.5, -1.7) {$[4, 6)$\quad 8};
  \node[s, minimum width=2.8cm, dim] at (4.5, -1.7) {$[6, 8)$\quad 6};
  \node[notel, align=center, text width=12.5cm] at (0, -2.6) {Adding 5 on $[2, 8)$: the two covering nodes (red) update their sums
    ($+5$ per element) and keep a tag; their ancestors are recomputed on the way back. The grey nodes are stale until some
    operation descends into $[4, 8)$ and pushes the tag down.};
\end{tikzpicture}
\caption{Lazy tags record updates that the subtree below still owes.}
\label{fig:lazy}
\end{figure}

\begin{fact}[What can be lazy]
Two things are needed: a tag must update a node's aggregate without looking at its children (``add $v$'' changes a sum by $v$ times the
length, a minimum by $v$), and two tags must \emph{compose} into one (``assign $x$, then add $y$'' is ``assign $x + y$''; ``add $y$,
then assign $x$'' is ``assign $x$''). Affine maps $a \mapsto ca + d$ compose, so ``multiply a range and add to a range'' is lazy; an
arithmetic progression added to a range (Polynomial Queries) is a tag $(\text{first term}, \text{step})$, and two of them add. ``Replace
each element by its square root'' does not compose into anything small; segment tree beats handle some such updates.
\end{fact}

\begin{problem}{Range Updates and Sums (CSES)}
An array receives range additions, range assignments and range-sum queries ($n, q \le 2 \cdot 10^5$).
\solution A lazy tree whose tag is the pair (assigned value or none, added amount), composed as in the fact above: 0.33\,s.
\end{problem}

\begin{code}
void solve() {  // LazyAssign (cplazyassign) is pasted above solve()
    int n, q;
    cin >> n >> q;
    vi a(n);
    for (ll& x : a) cin >> x;
    LazyAssign st(a);
    while (q--) {
        int type, l, r;
        cin >> type >> l >> r;
        l--;  // the range [l, r)
        if (type == 3) {
            print(st.sum_of(l, r));
            continue;
        }
        ll x;
        cin >> x;
        if (type == 1) st.add(l, r, x);
        else st.assign(l, r, x);
    }
}
\end{code}

\usetemplate{Lazy segment tree (add; sum, min) / (assign + add)}{Ctrl+Alt+T, Shift+S / Ctrl+Alt+Y, A}{cplazy, cplazyassign}

% ======================================================================================================================
\section[Counting minimums]{Counting minimums: the union of rectangles \hfill \plat}\label{sec:mincount}

A node that stores the minimum of its range \emph{and the total weight of the elements that attain it} supports range additions (the
minimum shifts, the weight does not change) and answers ``how much of the array is 0?'' whenever the values never go below 0. That is
exactly what a sweep over rectangles needs: sweep a vertical line from left to right, keep for every elementary $y$-interval how many
rectangles cover it, add $+1$ over a rectangle's $y$-range at its left edge and $-1$ at its right edge. Between two consecutive events
the covered length is the total length minus the length with count 0, and the area is that length times the width of the strip
(Figure~\ref{fig:rectsweep}).

\begin{figure}[htbp]
\centering
\begin{tikzpicture}[scale=0.62]
  \fill[fgsky!30] (0, 0) rectangle (4, 3);
  \fill[fgorange!25] (2, 1) rectangle (6, 5);
  \fill[fggreen!30] (2, 1) rectangle (4, 3);
  \draw[line width=0.9pt, draw=fgblue] (0, 0) rectangle (4, 3);
  \draw[line width=0.9pt, draw=fgorange!85!black] (2, 1) rectangle (6, 5);
  \foreach \xx in {0, 2, 4, 6} \draw[line width=0.8pt, draw=fgverm, dashed] (\xx, -0.5) -- (\xx, 5.5);
  \foreach \xx/\tx in {1/{$3 \cdot 2$}, 3/{$5 \cdot 2$}, 5/{$4 \cdot 2$}}
    \node[font=\sffamily\scriptsize\bfseries, text=fgverm, fill=white, inner sep=1pt] at (\xx, -0.95) {\tx};
  \draw[->, line width=0.5pt, fgink!70] (-0.6, 0) -- (7, 0) node[right, lbl] {$x$};
  \draw[->, line width=0.5pt, fgink!70] (0, -0.3) -- (0, 5.8) node[above, lbl] {$y$};
  \node[notel, anchor=west, text width=6.6cm] at (8.2, 2.5) {Events at $x = 0, 2, 4, 6$ (dashed). In the strip $[0, 2)$ the
    covered $y$-length is 3, in $[2, 4)$ it is 5 (the union $[0, 5)$, overlap counted once), in $[4, 6)$ it is 4. Area $= 6 + 10 + 8
    = 24 = 12 + 16 - 4$. The segment tree over the elementary $y$-intervals reports the uncovered length in $\Oh{1}$ after each
    event.};
\end{tikzpicture}
\caption{The union of two rectangles by a sweep line: the area is the sum of covered length times strip width.}
\label{fig:rectsweep}
\end{figure}

\begin{problem}{Area of Rectangles (CSES)}
Compute the area of the union of $n \le 10^5$ axis-parallel rectangles with coordinates up to $10^6$ in absolute value.
\end{problem}

\begin{code}
void solve() {  // MinCount (cpmincount) is pasted above solve()
    int n;
    cin >> n;
    vector<array<ll, 4>> ev;  // (x, +1 or -1, y1, y2)
    vi ys;
    rep(i, n) {
        ll x1, y1, x2, y2;
        cin >> x1 >> y1 >> x2 >> y2;
        ev.push_back({x1, 1, y1, y2});
        ev.push_back({x2, -1, y1, y2});
        ys.push_back(y1), ys.push_back(y2);
    }
    sort(all(ys));
    ys.erase(unique(all(ys)), ys.end());
    int m = sz(ys) - 1;  // elementary intervals [ys[j], ys[j + 1])
    vi len(m);
    rep(j, m) len[j] = ys[j + 1] - ys[j];
    MinCount st(len);  // how often each interval is covered, weighted
    ll total = ys.back() - ys[0], area = 0;
    sort(all(ev));
    rep(i, sz(ev)) {
        auto [x, d, y1, y2] = ev[i];
        int l = lower_bound(all(ys), y1) - ys.begin();
        int r = lower_bound(all(ys), y2) - ys.begin();
        st.add(l, r, d);
        if (i + 1 < sz(ev))  // the covered length times the strip's width
            area += (total - st.zero_weight()) * (ev[i + 1][0] - x);
    }
    print(area);
}
\end{code}

\usetemplate{Min + count segment tree}{Ctrl+Alt+Y, Shift+C}{cpmincount}

The weights are the lengths of the elementary intervals between consecutive distinct $y$-coordinates, so the tree has at most $2n$
leaves whatever the coordinates are; $10^5$ rectangles take 0.27\,s.

% ======================================================================================================================
\section[Dynamic and persistent segment trees]{Dynamic and persistent segment trees \hfill \plat}\label{sec:persistent}

\paragraph{Dynamic trees.} When positions go up to $10^9$ or $10^{18}$ and arrive online, so that they cannot be compressed in
advance, create nodes only when an update reaches them. Each update creates $\Oh{\log C}$ nodes for a coordinate range of size $C$; parts
of the range that nothing has touched do not exist and count as the identity.

\usetemplate{Dynamic segment tree}{Ctrl+Alt+Y, D}{cpdynseg}

\paragraph{Persistence.} A \defn{persistent} structure keeps every old version. An update in a segment tree changes one root-to-leaf
path, so copy just that path and let the new nodes point at the old, unchanged subtrees (Figure~\ref{fig:persistent}): $\Oh{\log n}$
new nodes per update, and every version stays queryable through its own root. Copying a whole array is free: the copy is just another
pointer to the same root.

\begin{figure}[htbp]
\centering
\begin{tikzpicture}
  \tikzset{pn/.style={vtxs, minimum size=6.5mm}}
  \node[pn] (r0) at (0, 0) {$v_0$};
  \node[pn] (a) at (-1.4, -1.1) {$a$}; \node[pn] (b) at (1.4, -1.1) {$b$};
  \node[pn] (c) at (0.7, -2.2) {$c$}; \node[pn] (d) at (2.1, -2.2) {$d$};
  \node[pn] (e) at (-2.1, -2.2) {}; \node[pn] (f) at (-0.7, -2.2) {};
  \draw[ed] (r0) -- (a) (r0) -- (b) (a) -- (e) (a) -- (f) (b) -- (c) (b) -- (d);
  \node[pn, hot] (r1) at (4.2, 0) {$v_1$};
  \node[pn, hot] (b1) at (5.4, -1.1) {$b'$};
  \node[pn, hot] (d1) at (6.1, -2.2) {$d'$};
  \draw[hded] (r1) -- (b1);
  \draw[hded] (b1) -- (d1);
  \draw[ded, draw=fgorange!90!black, dashed] (r1) to[bend right=12] (a);
  \draw[ded, draw=fgorange!90!black, dashed] (b1) to[bend left=20] (c);
  \node[notel, anchor=west, text width=5.0cm] at (7.0, -1.1) {Changing the leaf $d$ of version $v_0$ creates $v_1$: only the
    three red nodes $v_1$, $b'$ and $d'$ are new, and the dashed edges point into $v_0$'s unchanged subtrees $a$ and $c$. Both
    versions remain complete trees.};
\end{tikzpicture}
\caption{Path copying: a persistent update costs $\Oh{\log n}$ new nodes.}
\label{fig:persistent}
\end{figure}

\begin{problem}{Range Queries and Copies (CSES)}
Starting with one array, process: set a value in array $k$, ask for a range sum in array $k$, or append a copy of array $k$ as a new
array ($n, q \le 2 \cdot 10^5$).
\solution One persistent tree, one root per array. An assignment adds the difference along a copied path; a copy appends the root
pointer. Building the initial version directly costs $2n$ nodes instead of the $n \log n$ of $n$ insertions; the program runs in
0.29\,s.
\end{problem}

\begin{code}
void solve() {  // PersistentSeg (cppersist) is pasted above solve()
    int n, q;
    cin >> n >> q;
    vi a(n);
    for (ll& x : a) cin >> x;
    PersistentSeg ps(n);
    vi roots{-1, ps.build(a)};  // roots[k]: the current version of array k
    while (q--) {
        int type, k;
        cin >> type >> k;
        if (type == 3) {  // a copy shares every node with the original
            roots.push_back(roots[k]);
            continue;
        }
        int x;
        ll y;
        cin >> x >> y;
        if (type == 1) {  // a[x] = y: add the difference, copying one path
            ll cur = ps.query(roots[k], x - 1, x);
            roots[k] = ps.update(roots[k], x - 1, y - cur);
        } else print(ps.query(roots[k], x - 1, y));
    }
}
\end{code}

\usetemplate{Persistent segment tree}{Ctrl+Alt+Y, P}{cppersist}

\begin{trick}[The $k$-th smallest value in a range, online]
Build version $i$ of a persistent tree over the (compressed) values as ``the counts of $a_0, \dots, a_{i-1}$''. The difference of
versions $r$ and $l$ holds the counts of $a_l, \dots, a_{r-1}$, and a node-by-node difference can be descended without building it:
if the left children's counts differ by at least $k$, go left; otherwise subtract that difference and go right. That is $\Oh{\log n}$
per query, online, and the template's \cc{kth} does exactly this walk.
\end{trick}

% ======================================================================================================================
\section[Segment tree beats]{Segment tree beats \hfill \ioi}\label{sec:beats}

The update ``$a_i \leftarrow \min(a_i, x)$ for every $i$ in a range'' is not lazy in general: its effect on a sum depends on how many
elements exceed $x$. \defn{Segment tree beats} store in each node the maximum, the second largest distinct value and the number of
elements equal to the maximum. At a node, if $x \ge$ the maximum, nothing changes; if the second maximum $< x <$ the maximum, only the
maximal elements change, and the node updates its sum as $\mathit{sum} \mathrel{-}= (\mathit{max} - x) \cdot \mathit{count}$ and
keeps a tag; only otherwise does the update recurse. A potential argument (the number of distinct values in the nodes, which a
recursion strictly decreases) bounds the total work by $\Oh{(n + q) \log^2 n}$, and by $\Oh{(n + q) \log n}$ without range additions.
The same idea handles $\max$ updates and mixes of $\min$, $\max$ and additions.

\usetemplate{Segment tree beats}{Ctrl+Alt+Y, S}{cpbeats}

% ======================================================================================================================
\section[Offline sweeps]{Offline sweeps over queries \hfill \plat}\label{sec:sweepqueries}

Many range questions become one-dimensional when the queries are sorted and the array is \emph{swept}: a query $[l, r)$ is answered
at the moment the sweep reaches $r$, by a Fenwick tree over the positions or values to the left. The classic example is the number of
distinct values: sweep $r$ from left to right and keep a 1 at the \emph{last} occurrence of every value seen so far, moving it when a
value reappears. Then a value occurs in $[l, r)$ exactly when its last occurrence before $r$ is at least $l$, so the answer is the number
of 1s in $[l, r)$ (Figure~\ref{fig:sweepdistinct}).

\begin{figure}[htbp]
\centering
\begin{tikzpicture}
  \node[blbl, anchor=east] at (-0.2, 0) {$a$};
  \arrs{A}{0}{0}{1/n, 2/n, 1/n, 3/n, 2/n}
  \idx{A}{5}
  \node[blbl, anchor=east] at (-0.2, -1.35) {last occurrences};
  \arrs{B}{0}{-1.35}{0/dim, 0/dim, 1/good, 1/good, 1/good}
  \cellframe{B-2}{B-5}{fgverm}
  \node[notel, anchor=west, text width=7.4cm] at (4.0, -0.55) {After the sweep has passed 5 elements, the 1s mark the last
    occurrences of the values 1, 3 and 2. The query $[1, 5)$ counts the 1s in positions 1 to 4 (red frame): 3 distinct values,
    $\{2, 1, 3\}$. The first 1 and the first 2 were cleared when their values reappeared.};
\end{tikzpicture}
\caption{Distinct values by sweeping the right end: only the last occurrence of each value carries a 1.}
\label{fig:sweepdistinct}
\end{figure}

\begin{code}
void solve() {  // BIT (cpbit) is pasted above solve()
    int n, q;
    cin >> n >> q;
    vi a(n);
    for (ll& x : a) cin >> x;
    vector<vector<pii>> at(n + 1);  // the queries by right end: (l, index)
    rep(i, q) {
        int l, r;
        cin >> l >> r;
        at[r].push_back({l - 1, i});  // the range [l - 1, r)
    }
    // sweep r; the BIT holds a 1 at the last occurrence of every value in
    // a[0, r), so the 1s inside [l, r) are exactly the distinct values
    BIT b(n);
    map<ll, int> last;
    vi ans(q);
    rep(r, 1, n + 1) {
        ll v = a[r - 1];
        if (last.count(v)) b.add(last[v], -1);
        b.add(r - 1, 1), last[v] = r - 1;
        for (auto [l, i] : at[r]) ans[i] = b.sum(l, r);
    }
    for (ll x : ans) print(x);
}
\end{code}

This solves Distinct Values Queries in $\Oh{(n + q) \log n}$, 0.25\,s, the same as Mo's algorithm on these limits but without the
$\sqrt n$ factor. The library also has it as a function.

\usetemplate{Distinct values in ranges}{Ctrl+Alt+Y, Shift+D}{cpdistinct}

\paragraph{Containment as a sweep.} ``How many ranges does each range contain, and how many contain it?'' is a two-dimensional
condition: $[x_j, y_j] \subseteq [x_i, y_i]$ means $x_i \le x_j$ and $y_j \le y_i$. Sort the ranges by start, longer ones first among
equal starts; then every range \emph{before} $i$ starts no later than $i$, and the ones that contain $i$ are those that also end no
earlier, a count over the end points in a Fenwick tree. Sweeping backwards counts the contained ranges the same way. Identical ranges
contain each other and need one line of care: they are processed as a group.

\begin{problem}{Nested Ranges Count (CSES)}
For each of $n \le 2 \cdot 10^5$ ranges, count the ranges it contains and the ranges that contain it.
\end{problem}

\begin{code}
void solve() {  // BIT (cpbit) is pasted above solve()
    int n;
    cin >> n;
    vi x(n), y(n);
    rep(i, n) cin >> x[i] >> y[i];
    vi ys = y;
    sort(all(ys));
    ys.erase(unique(all(ys)), ys.end());
    auto rk = [&](ll v) { return lower_bound(all(ys), v) - ys.begin(); };
    vi ord(n);
    iota(all(ord), 0);
    // by start, longer first: every earlier range starts no later
    sort(all(ord), [&](int a, int b) {
        return pll{x[a], -y[a]} < pll{x[b], -y[b]};
    });
    auto same = [&](int a, int b) { return x[a] == x[b] && y[a] == y[b]; };
    vi inside(n), around(n);
    BIT bit(sz(ys));
    for (int s = 0, e; s < n; s = e) {  // one group of equal ranges at a time
        for (e = s; e < n && same(ord[e], ord[s]); e++) {}
        int i = ord[s];
        ll c = bit.sum(rk(y[i]), sz(ys)) + (e - s - 1);  // earlier, ends later
        rep(k, s, e) around[ord[k]] = c;
        bit.add(rk(y[i]), e - s);
    }
    bit = BIT(sz(ys));
    for (int s = n - 1, e; s >= 0; s = e) {  // backwards: every later range
        for (e = s; e >= 0 && same(ord[e], ord[s]); e--) {}  // starts later
        int i = ord[s];
        ll c = bit.prefix(rk(y[i]) + 1) + (s - e - 1);  // and ends sooner
        rep(k, e + 1, s + 1) inside[ord[k]] = c;
        bit.add(rk(y[i]), s - e);
    }
    print(inside);
    print(around);
}
\end{code}

\begin{recognize}
Offline queries whose answer depends on ``the last or first occurrence'', on the nearest greater element, or on pairs $(i, j)$ with
conditions in two coordinates: sort by one coordinate, sweep it, and keep the other in a Fenwick tree. Counting points in rectangles
the same way (each rectangle split into two prefix queries in $x$) takes $\Oh{(n + q) \log n}$ time and $\Oh{n}$ memory, far less than
a 2D Fenwick tree when the coordinates are large.
\end{recognize}

% ======================================================================================================================
\section[A staircase for dominance]{A staircase for two-dimensional dominance \hfill \ioi}\label{sec:staircase}

If pairs $(k, v)$ are only inserted, and a query asks for the maximum $v$ among the keys $\ge x$, a full segment tree is not needed. Keep
a \cc{map} from keys to values in which the values strictly \emph{decrease} as the keys increase: a new pair that is dominated (some key
at least as large already has a value at least as large) is not inserted, and otherwise it removes the pairs it dominates, which sit
just before it. The answer to a query is the value of the first key $\ge x$. Every pair is inserted and removed at most once, so the
operations are $\Oh{\log n}$ amortised. This staircase keeps two-dimensional dominance DPs, such as the longest chain of pairs
increasing in both coordinates, short and fast.

\begin{summary}[Chapter summary]
\begin{itemize}
  \item A Fenwick cell covers the $\mathit{lowbit}(i)$ positions ending at $i$: prefix sums remove lowest bits, updates add them. Two
        trees give range updates with range sums; descending by powers of two finds the $k$-th element in $\Oh{\log n}$.
  \item A segment tree query uses at most two nodes per level; a bottom-up array of $2n$ cells is the fastest implementation.
  \item Descending finds the first position where an aggregate reaches a value; a node stores whatever the merge of two halves needs.
  \item Lazy tags need a cheap effect on aggregates and closure under composition; min-count nodes measure unions of rectangles.
  \item Path copying makes a segment tree persistent with $\Oh{\log n}$ nodes per update; versions difference into $k$-th smallest
        queries.
  \item Sorting queries and sweeping one coordinate turns two-dimensional conditions into Fenwick-tree prefix sums.
\end{itemize}
\end{summary}

\begin{practice}
\ppgroup{Fenwick trees and segment trees}
\begin{pplist}
  \item Dynamic Range Sum Queries
  \item Dynamic Range Minimum Queries
  \item Range Update Queries
  \item Forest Queries II
  \item Hotel Queries
  \item List Removals
  \item Salary Queries
  \item Prefix Sum Queries
  \item Pizzeria Queries
\end{pplist}
\ppgroup{Nodes, lazy tags, persistence and sweeps}
\begin{pplist}
  \item Subarray Sum Queries
  \item Range Updates and Sums
  \item Polynomial Queries
  \item Range Queries and Copies
  \item Distinct Values Queries
  \item Nested Ranges Check
  \item Nested Ranges Count
  \item Area of Rectangles
  \item Increasing Array Queries
\end{pplist}
\end{practice}
